\section{Shared One-Window Processing and Streaming}

\subsection{The same detector for batch and edge use}

\path{WindowDetector.process()} is the common scoring and adaptation path.
Construction loads the backend and accepted detector once. One call returns
only the current event, threshold attempt, and refit attempt; absent work is
represented by \texttt{None}. Batch code collects those returns into complete
research arrays. Streaming consumes them immediately.

This arrangement avoids maintaining separate statistical rules for deployment.
It also preserves a useful distinction: the runtime is bounded, while a batch
research report intentionally grows with the finite experiment. The detector
retains bounded deques, active/original parameters, and counters rather than
accumulating every event or update attempt.

\subsection{The input and output interface}

The \texttt{run} command accepts UTF-8 JSON lines on standard input.
Each line has a device ID, a source ID, and a nested window with UTC time bounds,
ordered categories, counts, and observation state. The source ID becomes the
event's capture context; it need not contain a date or belong to a research
partition.

The maximum physical line is 65,536 bytes, including its ending. The reader
enforces that bound while reading, before a long line can enter memory.
It validates JSON, finite numeric values, required fields, device, count types,
category order, exact trained duration, and epoch-aligned UTC boundaries.
Timezone-aware offsets are normalised to UTC. Overlapping, duplicate, and
backward windows fail. A gap does not create an invented zero vector.

Missing windows advance chronology without producing events or learning data.
Observed silence is scored normally. A valid last record does not require a
newline. Blank, malformed, truncated, oversized, or incompatible records stop
with a line-numbered error and a nonzero exit code. Earlier emitted events
remain usable.

Standard output contains only the existing anomaly-event JSON, including raw
score, threshold, decision, severity, expected profile, residuals, and model
fingerprint. Each event is flushed before the next input read. Status and
accepted/rejected update reasons go to application logs; verbose mode also
uses standard error.

\subsection{Flow control and shutdown}

Input reading and event writing are synchronous. A slow consumer pauses further
reads instead of creating a growing output queue. Output and input errors
terminate the run. This keeps the memory design simple; it does not guarantee
delivery if a pipe or container log store fails.

SIGINT and SIGTERM stop an idle reader cleanly. During scoring or persistence,
the first stop request is deferred until the detector operation completes.
This protects atomic activation; it is not a hard guarantee that an arbitrary
future refit will finish within a service timeout. Signal handlers are restored
after the run, and the directory lock is released.

\subsection{Python access and tested equivalence}

The Python API accepts a validated \path{WindowRecord}; callers construct one
detector and hold the state lock around successive calls. The streaming
boundary adds line-size, chronology, and UTC-grid checks beyond the shared
per-window device/category/duration validation.

Twelve controlled CLI/batch combinations cover two score conventions, both
policies, and three modes. Comparisons include events, alpha, thresholds,
fingerprints, and accepted/rejected outcomes. Other tests verify the first
event arrives while input is still open, signal handling at an activation
boundary, idle shutdown, a slow output consumer, malformed inputs, state
restoration, and edge-configuration rollback.
